\subsection{Симметричные и бесконечные цепи}
%%%%%%%%%%%%%%%%%%%%%%%%%%%%%%%%%%%%%%%%%%%%%%%%%%%%
%%%%%%%%%%%%%%%%%%%%%%%%%%%%%%%%%%%%%%%%%%%%%%%%%%%%1
 \z[{\begin{wrapfigure}[4]{r}{0.2\textwidth}%
		\vspace{-25pt}%
		\begin{flushright}%
 			\includegraphics[width=0.18\textwidth]{3.8.1}%
		\end{flushright}%
	\end{wrapfigure}%
	}]%
{Определите токи в каждой стороне ячейки, полный ток от узла $A$ к узлу $B$ и полное сопротивление между этими узлами. Сторона каждой ячейки имеет сопротивление $r$, и ток, протекающий по одной из сторон, равен $i$.
\vspace{2em}
}
%
{$I = \frac{7}{2}i$; \quad $R = \frac{12}{7}r$}

%%%%%%%%%%%%%%%%%%%%%%%%%%%%%%%%%%%%%%%%%%%%%%%%%%%%1
 \z[{\begin{wrapfigure}[3]{r}{0.2\textwidth}%
		\vspace{-25pt}%
		\begin{flushright}%
 			\includegraphics[width=0.18\textwidth]{3.8.2}%
		\end{flushright}%
	\end{wrapfigure}%
	}]%
{Каждая сторона квадрата имеет сопротивление $r$. Определите сопротивление между узлами $A$ и $B$. Чему равно сопротивление между узлами $C$ и $D$?
\vspace{2em}
}
%
{$R_{AB} = \frac{13}{7}r$; \quad $R_{CD} = \frac{5}{7}r$}

%%%%%%%%%%%%%%%%%%%%%%%%%%%%%%%%%%%%%%%%%%%%%%%%%%%%1
 \z[{\begin{wrapfigure}[4]{r}{0.2\textwidth}%
		\vspace{-15pt}%
		\begin{flushright}%
 			\includegraphics[width=0.2\textwidth]{3.8.3}%
		\end{flushright}%
	\end{wrapfigure}%
	}]%
{Определите сопротивление $R_{AD}$ между точками $A$ и $D$ проволочной сетки, показанной на рисунке. Сопротивление каждого из проводников (вне зависимости от его длины), из которых спаяна сетка, равно $R$. Места спайки проводников обозначены точками. В центре сетки контакт отсутствует.
\vspace{2em}
}
%
{$R_{AD} = \frac{4}{5}R$}

%%%%%%%%%%%%%%%%%%%%%%%%%%%%%%%%%%%%%%%%%%%%%%%%%%%%1
 \z[{\begin{wrapfigure}[4]{r}{0.2\textwidth}%
		\vspace{-15pt}%
		\begin{flushright}%
 			\includegraphics[width=0.2\textwidth]{3.8.4}%
		\end{flushright}%
	\end{wrapfigure}%
	}]%
{Из проволоки сделали правильную четырёхугольную пирамиду, все рёбра которой имеют одинаковое сопротивление $R$. К серединам двух соседних (перпендикулярных) рёбер основания подсоединили выводы $A$ и $B$ омметра — прибора для измерения сопротивлений. Что покажет омметр?
\vspace{2em}
}
%
{$R_{AB} = \frac{2R}{3}$}

%%%%%%%%%%%%%%%%%%%%%%%%%%%%%%%%%%%%%%%%%%%%%%%%%%%%1
 \z[{\begin{wrapfigure}[3]{r}{0.2\textwidth}%
		\vspace{-30pt}%
		\begin{flushright}%
 			\includegraphics[width=0.2\textwidth]{3.8.5}%
		\end{flushright}%
	\end{wrapfigure}%
	}]%
{Определите эквивалентное сопротивление проволочной сетки между узлами $A$ и $B$, если сопротивление каждого проводника равно $R$.
\vspace{2em}
}
%
{$R_{AB} = \frac{11R}{19}$}

%%%%%%%%%%%%%%%%%%%%%%%%%%%%%%%%%%%%%%%%%%%%%%%%%%%%1
 \z[{\begin{wrapfigure}[1]{r}{0.3\textwidth}%
		\vspace{-20pt}%
		\begin{flushright}%
 			\includegraphics[width=0.3\textwidth]{3.8.6}%
		\end{flushright}%
	\end{wrapfigure}%
	}]%
{Найти сопротивление бесконечной цепочки.
\vspace{2em}
}
%
{$r = 2(1 + \sqrt{3})R$}

%%%%%%%%%%%%%%%%%%%%%%%%%%%%%%%%%%%%%%%%%%%%%%%%%%%%1
 \z[{\begin{wrapfigure}[4]{r}{0.4\textwidth}%
		\vspace{-0pt}%
		\begin{flushright}%
 			\includegraphics[width=0.4\textwidth]{3.8.7}%
		\end{flushright}%
	\end{wrapfigure}%
	}]%
{Какое сопротивление $r$ нужно присоединить к клеммам $C$ и $D$, чтобы сопротивление всей цепочки между клеммами $A$ и $B$ не зависело от числа элементарных ячеек?
}
%
{$r = (\sqrt{3} - 1)R$}


